\documentclass[10pt,a4paper]{amsart}
\usepackage[margin=19mm]{geometry}
\usepackage{amsmath,amssymb,mathtools}
\usepackage[scr=rsfs,cal=euler]{mathalfa}
\usepackage{xfrac}
\usepackage[unicode,hidelinks,bookmarks=false]{hyperref}
\newcommand{\ie}{i.e.\ }
\newcommand{\cf}{cf.\ }
\newcommand{\resp}{respectively\ }
\newcommand{\Subsection}{\subsection}
\providecommand{\nobreakdash}{\nobreak\hbox{-}}
\setlength{\emergencystretch}{2em}
\multlinegap0pt
\usepackage{amssymb}
\usepackage{xcolor}
\usepackage{subfig}
\usepackage{multirow}
\usepackage{booktabs}
\usepackage{breqn}
\usepackage{bm}
\usepackage[normalem]{ulem}
\newtheorem{proposition}{Proposition}
\title{Nonparametric Kernel Regression for Coordinated Energy Storage Peak Shaving with Stacked Services}
\author{Emily Logan, Ning Qi, Bolun Xu}
\date{Source: arXiv:2602.16586v1 (CC BY 4.0)}
\begin{document}
\maketitle

\noindent\emph{Source and licence.} Nonparametric Kernel Regression for Coordinated Energy Storage Peak Shaving with Stacked Services. Version arXiv:2602.16586v1 (CC BY 4.0), distributed by arXiv under the Creative Commons Attribution 4.0 International licence, http://creativecommons.org/licenses/by/4.0/, as recorded in the arXiv licence metadata for this exact version and shown on the abstract page https://arxiv.org/abs/2602.16586v1. Full licence notice: \texttt{licenses/arxiv-2602.16586-CC-BY-4.0.txt}.\par\smallskip

\noindent\textit{Editorial scope.} Counted unit: the article's proposition Equivalence\_Prop (source0.tex lines 205--207). The article's complete original proof of it is delivered with it (source0.tex lines 209--242, one \texttt{proof} environment of 34 lines). No mathematics is written, repaired or re-derived: the statement and the proof are the article's own lines, in the article's own notation, and no retained line is substituted or re-flowed.  Retained uncounted as the source-local context the proof needs: lines 127--142; lines 155--161; lines 165--180.  Nothing was omitted from any retained span, so the package carries no omission record.  No retained line contains a citation, so the wrapper carries no bibliography and no bibliography entry.  No retained line contains a figure environment, an image-inclusion command or drawing code, so no image is delivered and the package declares no asset dependency.  Editorial formatting only, in addition: an AMS \texttt{amsart} preamble that restates the article's own declarations and macros used by the retained lines, namely the environments \texttt{proposition} on separate counters, together with the standard packages those lines need.  The retained environments print in this document as Proposition 1 in order of appearance, whereas the article prints the same items with the numbering of its own full document.  The complete native source of the article as distributed in the arXiv e-print for this exact version is bundled unchanged as \texttt{sources/194902/source0.tex}.
\begin{subequations}\label{PS_and_Arb_Formulation}
\begin{align}
F_{\text{comb}}(x) &= \min_{p,d,q} \sum\nolimits_{t=1}^T 
\left[ c d_t + \lambda_t (D_t - d_t + q_t) \right] + \kappa p
\label{det_obj} \\
\text{s.t. } \quad
& p \ge D_t - d_t + q_t, \quad \forall t = 1, \dots, T
\label{peak_def} \\
& e_{t+1} = e_{t} - {d_t}/{\eta} + \eta q_t, \quad \forall t = 1, \dots, T
\label{soc_update} \\
& 0 \le d_t, q_t \le \overline{P}, \quad \forall t = 1, \dots, T
\label{power_bounds} \\
& \underline{E} \le e_{t+1} \le \overline{E}, \quad \forall t = 1, \dots, T
\label{energy_bounds}
\end{align}
\end{subequations}

\begin{equation}\label{PS_Only_Formulation}
\begin{aligned}
& F_{\text{PS}}(x;\delta) = \min_{p,d,q,e} \; p + \delta \sum\nolimits_{t=1}^T e_{t+1} \\
\text{s.t. } \quad
& \text{Constraints \eqref{peak_def}--\eqref{energy_bounds}}
\end{aligned}
\end{equation}

\textbf{Arbitrage.} 
Given the Stage 1 solution ($p^*$, $e^*$), Stage 2 maximizes the revenue from energy arbitrage while ensuring that arbitrage decisions respect both the peak shaving constraint $p \le p^*$ and the SoC trajectory established in Stage 1.
\begin{subequations}\label{Arb_Formulation_Stage2}
\begin{align}
& F_{\text{arb}}(x) = \min_{d,q,e} \sum\nolimits_{t=1}^T \big[ c d_t - \lambda_t (d_t - q_t) \big]
\label{arb_obj_min} \\
\text{s.t. } \quad
& p^* \ge D_t - d_t + q_t, 
\quad \forall t = 1,\dots,T
\label{peak_const_stage2} \\
& e_{t+1} \ge e_{t+1}^*, 
\quad \forall t = 1,\dots,T
\label{PS_SOC_min_stage2} \\
&  \eqref{soc_update}-\eqref{energy_bounds}. \nonumber
\end{align}
\end{subequations}

\begin{proposition}\label{Equivalence_Prop}
    Given that $\delta$ is sufficiently small and $\kappa \gg \lambda_t, c$, solving the two-stage optimization in \eqref{PS_Only_Formulation} and \eqref{Arb_Formulation_Stage2} is equivalent to solving the combined peak shaving and arbitrage problem in \eqref{PS_and_Arb_Formulation}.
\end{proposition}

\begin{proof}
We first restate the combined problem for convenience:
\begin{equation}
    F_{\text{comb}}(x)=\sum\nolimits_{t=1}^T \lambda_t D_t + F_{\text{arb}}(x) + \kappa F_{\text{PS}}(x). \nonumber
\end{equation}
Let $x=(p,\{d_t,q_t,e_{t+1}\}_{t=1}^T)$ denote the decision variables, and let $S$ be the feasible region defined by \eqref{peak_def}–\eqref{energy_bounds}, where $g(x)\le0$ collects all inequality constraints and $h(x)=0$ the equality constraint.  

The Lagrangian of the combined problem is
\begin{equation}
    \mathcal{L}_{\text{comb}}(x,\mu,\nu)=F_{\text{arb}}(x)+\kappa F_{\text{PS}}(x)+\mu^T g(x)+\nu^T h(x),
\end{equation}
with $\mu\ge0$ and $\nu$ denoting the dual variables. The KKT stationarity condition is
\begin{equation}\label{eq:kkt_combined}
    \nabla_x F_{\text{arb}}(x)+\kappa\nabla_x F_{\text{PS}}(x)+\nabla g(x)^T\mu+\nabla h(x)^T\nu=0,
\end{equation}
accompanied by primal and dual feasibility and complementary slackness.  

To analyze the limit $\kappa \gg \lambda_t, c$, define rescaled dual variables $\tilde{\mu}=\mu/\kappa$ and $\tilde{\nu}=\nu/\kappa$. Substituting into \eqref{eq:kkt_combined} and dividing by $\kappa$ yields
\begin{equation}\label{eq:kkt_rescaled}
   \nabla_x F_{\text{arb}}(x)/\kappa+\nabla_x F_{\text{PS}}(x)+\nabla g(x)^T\tilde{\mu}+\nabla h(x)^T\tilde{\nu}=0.
\end{equation}
Taking the limit as $\kappa\to\infty$ removes the first term, giving
\begin{equation}\label{eq:kkt_ps_stage1}
    \nabla_x F_{\text{PS}}(x)+\nabla g(x)^T\tilde{\mu}+\nabla h(x)^T\tilde{\nu}=0,
\end{equation}
which are precisely the KKT conditions of the peak shaving problem. Thus, for $\kappa\gg\lambda_t, c$ and $\delta\to0^+$, the solution to \eqref{PS_and_Arb_Formulation} coincides with the peak shaving optimum: $p_{\text{comb}}^*=p^*$ and $e_{\text{comb}}^*=e^*$.  

Since the combined solution minimizes $F_{\text{PS}}(x)$ under these conditions, it lies in the optimal peak shaving set $S^*=\{x\in S\mid F_{\text{PS}}(x)=F_{\text{PS}}^*\}$. Restricting the combined problem to $S^*$, we have:
\[
    \min_{x\in S^*}\big[F_{\text{arb}}(x)+\kappa F_{\text{PS}}(x)\big].
\]
As $F_{\text{PS}}(x)=F_{\text{PS}}^*$ for all $x\in S^*$, the constant term $\kappa F_{\text{PS}}^*$ can be dropped, leaving $\min_{x\in S^*}F_{\text{arb}}(x)$, identical to the arbitrage formulation \eqref{Arb_Formulation_Stage2} over the feasible set characterized by $p\le p^*$ and $e_{t+1}\ge e_{t+1}^*$ for all $t$.  

Therefore, when $\kappa$ is sufficiently large, solving the combined problem \eqref{PS_and_Arb_Formulation} is equivalent to the two-stage procedure of first solving the peak shaving problem \eqref{PS_Only_Formulation} to obtain $(p^*,e^*)$, and then optimizing arbitrage over $S^*$.  \end{proof}

\end{document}
